%% kapa: exam=1 round=32 year=2021 period=1 form=A subject=civil_law subject_ko=민법 seq=24 number=24
%% source: Q-Net 감정평가사 기출문제 (cst00309) artlSeq=5208397 file=extracted/1008/5208397/A형/2021년도 제32회 감정평가사 1차 1교시 A형.pdf
%% conversion: kordoc 4.24.51 -> tools/gen1.py (auto); review=proofread-sonnet; flags=none
\begin{question}{24}
점유에 관한 설명으로 옳지 않은 것은? (다툼이 있으면 판례에 따름)
\begin{choices}
  \item 점유매개자의 점유를 통한 간접점유에 의해서도 점유에 의한 시효취득이 가능하다.
  \item 사기의 의사표시에 의해 건물을 명도해 준 자는 점유회수의 소권을 행사할 수 없다.
  \item 미등기건물을 양수하여 건물에 관한 사실상의 처분권을 보유한 양수인은 그 건물부지의 점유자이다.
  \item 간접점유의 요건이 되는 점유매개관계는 법률행위가 아닌 법령의 규정에 의해서는 설정될 수 없다.
  \item 상속에 의하여 점유권을 취득한 상속인은 새로운 권원에 의하여 자기 고유의 점유를 개시하지 않는 한 피상속인의 점유를 떠나 자기만의 점유를 주장할 수 없다.
\end{choices}
\end{question}
\answer{4}
